% Vietnamese translation for OI Public Library
% Source-File: IOI中国国家候选队论文/2007/day2/5.陈雪《问题中的变与不变》.ppt
\documentclass[11pt]{article}
\usepackage{oipl-vn}

\newcommand{\Oh}{\mathcal{O}}
\newcommand{\cost}{\operatorname{cost}}

\title{Bản trình chiếu: Biến đổi và bất biến trong bài toán}
\author{Chen Xue}
\date{2007}

\begin{document}
\maketitle

\origtitle{Biến đổi và bất biến trong bài toán}
\sourcepath{IOI China candidate team papers / 2007 / day 2 / Chen Xue.ppt}

\section{Dẫn nhập}

Trong thi tin học, nhiều bài toán thực chất là quá trình tìm, cập nhật hoặc
duy trì các biến. Biến có thể là biến vòng lặp, biến tích lũy, biến quyết
định, hoặc những thông tin được lưu trong cấu trúc dữ liệu như cây cân bằng,
ngăn xếp và hàng đợi.

Khi quy mô lớn, chỉ riêng việc duy trì mọi biến thay đổi trực tiếp đã có thể
vượt giới hạn thời gian và bộ nhớ. Tư tưởng của slide là tìm phần ``không
đổi'' ẩn trong các biến: gom nhiều biến thành tập, tổng, tích, quan hệ thứ
tự, hoặc một biểu thức ít thay đổi hơn. Khi tìm đúng đại lượng bất biến, bài
toán mô phỏng có thể chuyển thành bài toán tính trực tiếp.

\section{Ví dụ 1: Ants}

Có nhiều con kiến bò trên đoạn thẳng dài \(l\) cm với vận tốc \(1\) cm/s. Khi
đến đầu mút thì kiến rơi xuống. Khi hai con kiến gặp nhau, chúng lập tức quay
đầu. Biết \(l\) và vị trí ban đầu của từng con kiến, nhưng không biết hướng
đi ban đầu. Hỏi thời điểm sớm nhất và muộn nhất mà tất cả kiến đều rơi xuống.
Số kiến có thể tới \(1\,000\,000\).

Nếu liệt kê hướng ban đầu, có \(2^n\) trường hợp. Nếu đã biết hướng rồi mà
vẫn mô phỏng va chạm, số lần gặp nhau có thể đạt \(\Oh(n^2)\). Biến khó duy
trì nhất là hướng của từng con kiến sau mỗi lần gặp.

Xét hai con kiến \(A,B\) gặp nhau tại vị trí \(X\). Ngay trước va chạm, tập
thông tin vận tốc-vị trí của chúng là
\[
  U=\{(1,X),(-1,X)\}.
\]
Ngay sau va chạm,
\[
  U'=\{(-1,X),(1,X)\}.
\]
Hai tập bằng nhau. Nếu không gắn danh tính cho kiến, sự kiện ``quay đầu''
tương đương với việc hai con kiến đi xuyên qua nhau và đổi tên. Vì vậy mọi
va chạm trong tập kiến không làm thay đổi tập các cặp vận tốc-vị trí.

Do đó tập thời điểm rơi xuống chính là
\[
  T=\{\text{thời gian để kiến }i\text{ đi theo hướng ban đầu tới một đầu mút}\}.
\]
Để thời điểm tất cả rơi là sớm nhất, cho mỗi kiến đi về đầu mút gần hơn, rồi
lấy giá trị lớn nhất. Để thời điểm muộn nhất, cho mỗi kiến đi về đầu mút xa
hơn. Thuật toán chỉ cần duyệt các vị trí một lần, đạt \(\Oh(n)\), cũng là
cận dưới vì phải đọc đầu vào.

\section{Bài học từ ví dụ kiến}

Điểm mấu chốt là bỏ danh tính cá thể. Hướng của từng con kiến thay đổi liên
tục, nhưng tập các chuyển động không đổi dưới phép va chạm. Đây là một mẫu
rất thường gặp: nếu phép biến đổi chỉ hoán đổi nhãn, ta nên mô tả hệ bằng
tập hoặc đa tập thay vì mô tả từng phần tử có tên.

Các đại lượng bất biến có thể là parity, tổng, tích, đa tập, thứ tự tương
đối, hoặc giá trị của một hàm trên trạng thái. Câu hỏi thực dụng khi đọc đề
là: sau mỗi thao tác, điều gì thật sự thay đổi, và điều gì chỉ đổi cách biểu
diễn?

\section{Ví dụ 2: Navigation Game}

Bài toán cho một lưới \(N\times M\). Từ một bến xuất phát ở hàng cuối, con
thuyền cần đi tới một bến ở hàng đầu. Thuyền chỉ đi trái, phải hoặc lên; không
được đi vào đất liền trừ khi cập bến; đã rời một ô thì không quay lại. Đi lên
tốn một đơn vị thời gian. Nếu trước đó đã đi ngang liên tiếp \(x\) lần, bước
ngang tiếp theo tốn \(x+1\). Trên biển còn có chướng ngại \(O\), ô số phận
\(F\) phải đi qua lẻ lần, đá chúc phúc \(B\) đi vào không tốn thời gian, và
bùa bão \(S\) làm chi phí bình thường bị nhân đôi. Giới hạn là
\(N,M\le 1000\).

Trạng thái DP cơ bản là \(f_{t,i,j}\): chi phí nhỏ nhất để đến ô \((i,j)\)
với parity số ô \(F\) đã đi qua là \(t\). Khi chuyển từ hàng \(i+1\) lên hàng
\(i\), chọn một cột \(k\) rồi đi ngang trong hàng \(i\):
\[
  f_{t,i,j}=
  \min_k\{f_{t',i+1,k}+C_{i,k}+\cost(i,k,j)\}.
\]
Nếu không có \(B,S\), chi phí ngang từ \(k\) tới \(j\) là
\[
  \cost(i,k,j)=\frac{|j-k|(|j-k|+1)}2,
\]
và chuyển trạng thái trở thành dạng tối ưu hóa độ dốc quen thuộc. Khi chỉ xét
từ trái sang phải, đặt \(U_j=f_{t,i,j}\), \(V_k=f_{t',i+1,k}+C_{i,k}\), ta có
\[
  U_j=\min_k\left\{V_k+\frac{(j-k)(j-k+1)}2\right\}.
\]
So sánh hai quyết định \(k,l\) đưa tới một điều kiện về hệ số góc của các
điểm
\[
  X_k=2V_k+k(k-1),\qquad Y_k=2k.
\]
Vì vậy có thể duy trì bao lồi dưới bằng hàng đợi và trả lời các quyết định
tối ưu trong thời gian khấu hao hằng số.

\section{Giữ bất biến trong Navigation Game}

Khó khăn là các ô \(B\) và \(S\). Chi phí \(\cost(i,k,j)\) không còn chỉ là
tam giác đơn giản, vì mỗi bước ngang có thể bị xóa chi phí hoặc nhân đôi. Tuy
nhiên với \(k\) cố định và \(j\) tăng, nó vẫn có dạng
\[
  \cost(i,k,j)=\frac{(j-k+1)(j-k)}2+a(j-k)+b.
\]
Khi ô mới không đặc biệt, \(a,b\) cập nhật theo một cách; khi là \(B\), hệ số
\(a\) giảm; khi là \(S\), hệ số \(a\) tăng. Ý nghĩa của \(a\) là tổng ảnh
hưởng của các ô \(S\) trừ các ô \(B\) trên đoạn đang xét.

Thay vì cố duy trì riêng mọi \(\cost(i,k,j)\), slide điều chỉnh biến
\(g_{t',i,k}\) để hấp thụ phần \(b\). Khi đó với một trạng thái \(f_{t,i,j}\),
các hệ số \(a,b\) được xem như cố định, và bất đẳng thức so sánh hai quyết
định vẫn có dạng độ dốc:
\[
 \frac{(2g_k+k(k-1))-(2g_l+l(l-1))}{2(k-l)}
 < j-\frac a2.
\]
Mỗi lần \(j\) tăng, vế phải vẫn không giảm, vì \(a\) chỉ thay đổi một đơn vị
theo ô \(B/S\). Do đó hàng đợi bao lồi vẫn hợp lệ. Làm tương tự cho hướng
phải sang trái, tổng độ phức tạp trở thành \(\Oh(NM)\).

Thông điệp của ví dụ này là không phải lúc nào đại lượng bất biến cũng hiện
ra dưới dạng ``không đổi tuyệt đối''. Đôi khi ta cần chuyển một phần biến
đổi sang biến khác để giữ nguyên hình dạng toán học của chuyển trạng thái.

\section{Ví dụ 3: Circular Railway}

Trên đường ray tròn dài \(L\) có \(n\) điểm loại \(A\) và \(n\) điểm loại
\(B\). Cần ghép mỗi điểm \(A\) với một điểm \(B\) sao cho tổng khoảng cách
ngắn nhất trên vòng tròn là nhỏ nhất. Giới hạn \(n\le 50000\), nên không thể
dùng ghép cặp trọng số tổng quát.

Sắp \(A_1,\ldots,A_n\) và \(B_1,\ldots,B_n\) theo chiều kim đồng hồ, đồng
thời nhân đôi dãy \(B\) để xử lý vòng tròn. Tính chất quan trọng là trong
ghép tối ưu, hai cạnh ghép không cắt nhau. Nếu có hai cạnh \(A_i-B_v\) và
\(A_j-B_u\) cắt nhau, đổi thành \(A_i-B_u\) và \(A_j-B_v\) không làm tổng
chi phí tăng. Vì vậy nếu \(A_1\) ghép với \(B_k\), thì \(A_i\) sẽ ghép với
\(B_{k+i-1}\).

Thuật toán \(\Oh(n^2)\) là thử mọi \(k\) và tính tổng
\[
  \sum_i C_i,\qquad C_i=\operatorname{dist}(A_i,B_{k+i-1}).
\]
Nhưng trực tiếp tính mọi \(C_i\) vẫn quá đắt. Cần tìm đại lượng ít thay đổi
khi \(k\) tăng.

\section{Các ``hằng lượng'' trong Circular Railway}

Khoảng cách ngắn nhất giữa \(A_i\) và \(B_j\) thuộc một trong bốn dạng:
\[
  A_i-B_j,\qquad B_j-A_i,\qquad L+A_i-B_j,\qquad L+B_j-A_i,
\]
tùy vị trí tương đối và việc đi theo cung ngắn hơn. Mỗi dạng có thể viết
thành
\[
  C_i=f(A_i)+g(B_j).
\]
Khi \(k\) tăng từ \(1\) đến \(n\), với mỗi \(A_i\), giá trị \(f(A_i)\) chỉ
đổi ở vài ngưỡng nơi \(B_j\) đi qua \(A_i\) hoặc đi qua điểm cách nửa vòng.
Tương tự \(g(B_j)\) cũng chỉ đổi vài lần. Vì vậy từng hạng tử không còn được
xem là biến phức tạp; nó được mô tả bằng một số hữu hạn sự kiện.

Tiền xử lý các thời điểm sự kiện bằng hai con trỏ trên hai dãy đã sắp. Khi
duyệt \(k=1..n\), ta cập nhật tổng hiện tại bằng các sự kiện xảy ra ở thời
điểm \(k\), rồi lấy giá trị nhỏ nhất. Tổng số sự kiện không quá hằng số nhân
\(n\), nên độ phức tạp là
\[
  \Oh(n\log n+n),
\]
trong đó phần \(\Oh(n\log n)\) đến từ sắp xếp.

\section{Kết luận}

Ba ví dụ minh họa ba kiểu tìm bất biến:
\begin{itemize}
  \item Ants: bỏ danh tính cá thể để thấy va chạm chỉ hoán đổi nhãn;
  \item Navigation Game: đổi biến để giữ nguyên dạng chuyển trạng thái tối ưu
  hóa độ dốc;
  \item Circular Railway: tách chi phí thành các phần chỉ đổi tại một số ít
  sự kiện.
\end{itemize}
Kỹ thuật chung là quan sát kỹ thao tác biến đổi, mạnh dạn đoán đại lượng thật
sự cần duy trì, rồi chứng minh đại lượng ấy đủ để suy ra đáp án.

\end{document}
